% Template: 한국정보과학회 학술대회 논문작성 양식 (latex)
% Main sources
% - http://ids.snu.ac.kr/wiki/LaTeX_template_for_KCC
% - https://github.com/leejaymin/Kiise-Latex-template
% Edited by Jinhong Jung (Soongsil University, DMLAB)

\documentclass[preprint]{kcc} % for review
%\documentclass{kcc} % for publication

% 기본 패키지 외에 표나 코드 삽입을 위해 필요한 패키지 추가
\usepackage{booktabs}
\usepackage{listings}
\usepackage{graphicx}
\usepackage{multirow}

%\input{packages}
%\input{definitions}
							
% set title, author, abstract
\title{cuTile 기반 고성능 LLM 추론 엔진 최적화 및 코딩 에이전트 협력 설계 실증 연구}
\author{
저자1 국문이름$^{\circ}$, 저자2 국문이름\\ % HTML 원문에 개별 저자명은 없으므로 양식 유지
우송대학교 소프트웨어학부(AI·컴퓨터공학과) \\
\{author1email, author2email\}@wsu.ac.kr
}
\engtitle{Optimizing High-Performance LLM Inference Engine via cuTile and Human-Agent Co-Design: An Empirical Study}
\engauthor{
First Author Name, Second Author Name\\
Dept. of AI and Computer Engineering, Woosong University\\
}
\begin{document}
\abstract{
거대 언어 모델(LLM) 서빙 엔진의 지연 시간 및 처리량 성능은 GPU 메모리와 연산 대역폭 병목을 해결하기 위한 하드웨어 수준의 커널 최적화에 직결된다. 그러나 vLLM 등 대다수의 고성능 서빙 프레임워크는 Linux 환경에 강력하게 종속되어 있어, 연구 및 교육 목적과 로컬 온프레미스 장비가 집중된 Windows OS 생태계로의 이식성에 심각한 장벽이 존재한다. 본 연구는 Windows OS 환경에서 NVIDIA의 최신 JIT 컴파일 기반 Python DSL인 cuTile 및 CUDA Python 1.0을 도입하여 크로스 플랫폼 호환성을 보장하는 고성능 추론 엔진인 \textbf{micro-vllm}을 설계하고 구현하였다.
특히, 복잡한 GPU 커널 코딩과 E2E 시스템 설계 병목의 탐색 공간을 효율적으로 최적화하기 위해, 별도의 자율형 에이전트 루프 시스템 구축 대신 \textbf{AI 코딩 에이전트(Coding Agent)}를 공동 설계(Co-design) 및 병목 분석 파트너로 활용하는 인간-에이전트 협력적 최적화 파이프라인을 구축하였다. 본 논문은 이 파이프라인을 바탕으로 미시적 커널 수준에서 레지스터 스필링(Register Spilling)을 유발하는 물리적 하드웨어 임계치를 규명하고 64x64 타일링 스케줄링을 도출하였으며, Paged KV 캐시 상의 dynamic shape mismatch를 sub-millisecond 수준에서 해소하는 벡터화된 논리-물리 블록 재배치 기법을 고안하였다.
나아가 시스템 이식 과정에서 JIT 컴파일 방지를 위한 dynamic padding 기법이 오히려 PyTorch caching allocator와 가비지 컬렉터의 메모리 할당 스래싱을 촉발하여 전체 생성 처리량을 67.0\% 저하시키는 시스템적 부작용을 규명하고 무패딩 Eager Serving 구조를 확립하였다. 최종적으로 CUDA Green Contexts 기술을 적용하여 SM 자원을 Prefill(32 SM)과 Decode(16 SM) 영역으로 물리 격리함으로써 단일 소비자용 GPU(RTX 5070) 환경에서 \textbf{중간값 Decode 지연시간(P50 ITL) 14.0\% 단축 및 전체 토큰 처리량 10.6\% 향상(459.73 tok/s)}을 달성하였다.

\par\smallskip\noindent\textbf{핵심어(Keywords):} 거대 언어 모델(LLM) 서빙, cuTile, CUDA Python 1.0, Windows OS, 코딩 에이전트 공동 설계, Green Contexts, 메모리 스래싱
}



\maketitle

\section{서론}
최근 생성형 AI 기술의 급격한 성장에 힘입어 로컬 온프레미스 장비나 오프라인 연구실 환경에서 LLM을 구동하려는 수요가 폭증하고 있다. 특히 RAG(Retrieval-Augmented Generation), 로컬 지식베이스 구축, 그리고 기업 정보 보안을 유지하기 위한 로컬 프라이빗 추론 환경 구축은 현업 시스템 설계에 있어 필수적인 과제이다. 그러나 vLLM이나 TensorRT-LLM 등 실용적으로 널리 쓰이는 대규모 서빙 아키텍처들은 Triton 커널 컴파일러, FlashAttention 패키지, 그리고 분산 텐서 통신용 NCCL 라이브러리 등 강력한 Linux 독점 환경에 강하게 묶여 있다. 이로 인해 데스크톱 컴퓨터와 로컬 서빙의 보편적 운영체제인 Windows 환경에서는 WSL2(Windows Subsystem for Linux) 우회를 통한 가상화 비용과 물리 기기 다이렉트 제어 제약 등으로 최상의 연산 성능을 보장하기 어려운 한계가 존재했다.

본 연구는 이러한 크로스 플랫폼 장벽을 해소하고 최적의 로컬 LLM 서빙 환경을 제공하기 위해, Windows OS 환경에서 동작하는 NVIDIA의 새로운 \textbf{cuTile} Python DSL 및 \textbf{CUDA Python 1.0} 기반의 고성능 서빙 엔진인 \textbf{micro-vllm}을 개발하였다. 이 과정에서 C++ 수준의 직접적인 CUDA 프로그래밍 오버헤드와 하드웨어 프로파일링 분석의 긴 탐색 주기를 단축시키기 위해, \textbf{AI 코딩 에이전트(Coding Agent)}와의 협력적 튜닝 파이프라인을 가동하여 시스템을 설계, 실험 및 개선해 나가는 실험과학적 엔지니어링 접근을 도입하였다.

본 논문은 단일 연산 커널 수준(Micro-level)에서의 최적화 구조가 실제 엔드투엔드(E2E) 서비스 아키텍처로 통합될 때 야기되는 '추상화의 누수(Abstraction Leak)' 병목 현상을 실증적으로 추적하고 그 대안을 제시한다. 구체적으로는 dynamic padding 최적화가 유발하는 PyTorch CUDA caching allocator의 호스트-디바이스 메모리 스래싱 문제를 규명하였으며, CUDA Green Contexts 기술을 통한 하드웨어 자원 격리 벤치마크 결과를 분석하여 디코딩의 메모리 대역폭 향상과 캐시 효율성을 입증하였다.

\subsection{본 연구의 주요 기여}
본 논문이 제시하는 학술적 및 시스템적 주요 기여점은 그 중요도 순에 따라 다음과 같다.

\begin{itemize}
    \item \textbf{[Core Discovery] 역직관적 시스템 병목 규명 (Allocator Thrashing)}: JIT 방지를 위해 형태를 패딩하는 기존 상식과 직관을 뒤집고, Eager-mode LLM serving 환경 하에서 dynamic padding이 오히려 메모리 할당자의 catastrophic contention(이중 연결 리스트 스캔, cudaMalloc/cudaFree, GC 동기화 오버헤드 등)을 촉발하여 전체 처리량을 67.0\%나 붕괴시킨다는 역설적인 아키텍처적 병목을 규명하였다.
    \item \textbf{[AI4Sys] 지식-주도형 코딩 에이전트 협력 프레임워크 (Human-AI Co-design Framework)}: 하드웨어 프로파일링의 물리적 한계점(예: register spilling)을 영구적인 DSL 제약 규칙으로 동적 피드백하여 컴파일 성능을 보증하는 3계층(Architect-TechLead-Executor) 협업 최적화 방법론을 제시하였다. 이를 통해 복잡한 GPU 커널 개발이 자율 진화 순환형(Evolutionary Closed Loop) 과학실험 프레임워크로 통제되어 가속될 수 있음을 증명하였다.
    \item \textbf{[System \& Methodology] Windows-네이티브 구현 및 자원 격리 솔루션 (Windows-Native serving with Green Contexts)}: 리눅스 의존성이 심각한 서빙 환경을 극복하고 Windows 환경에서 이식 및 가속을 보증하기 위해, cuTile 및 CUDA Python 1.0 기반 고성능 엔진인 \textbf{micro-vllm}을 설계하고 실증적으로 구현했다. 특히 CUDA Green Contexts 기술로 SM 자원을 Prefill(32 SM)과 Decode(16 SM)로 강제 구획 분할하여 L2 캐시 residency를 격리함으로써, Decode P50 ITL을 14.0\% 단축하고 Throughput을 10.6\% 개선하는 물리 격리 솔루션을 정량 입증했다.
    \item \textbf{[Educational \& Reproducibility] 투명한 단계별 커널 커리큘럼 (Progressive serving stack without black-box libraries)}: 기존의 거대 서빙 라이브러리들이 의존하는 블랙박스 커널(FlashAttention 등) 종속성에서 탈피하여, 단일 GEMM부터 FMHA, Paged KV 캐시 레이어에 이르기까지 전체 추론 스택의 핵심 요소를 cuTile로 progressive하게 구현하였다. 이를 통해 GPU 아키텍처 교육과 실제 운영 등급의 시스템 설계 간의 지식 격차를 메우고 학계 및 교육 환경에서의 완전한 재현성(Reproducibility)을 제공한다.
\end{itemize}


\section{관련 연구 및 배경}
\subsection{LLM 서빙 엔진과 PagedAttention}
LLM의 자기회귀적(Autoregressive) 생성 루프는 매 스텝마다 과거의 토큰 정보(Key, Value)를 저장하고 재사용하는 KV 캐싱을 핵심 구조로 가진다. vLLM은 이 KV 캐시를 고정 크기의 물리적 페이지 블록으로 분할 관리하는 PagedAttention 기법을 도입하여 GPU 메모리의 외부/내부 단편화를 해결하였으며, 기존의 메모리 관리 기법과 달리 처리량을 대폭 향상시켰다. 그러나 기존 시스템들은 분산 환경을 전제로 NCCL 통신 레이어에 의존하고 있어, 단일 혹은 소비자용 GPU 환경 및 Windows와 같은 이기종 OS에서의 호환성이 현저히 떨어진다.

\subsection{CUDA Python 1.0 및 cuTile DSL}
최근 공식 릴리즈된 CUDA Python 1.0은 Python 런타임 환경에서 하드웨어 드라이버 및 CUDA API를 직접 조작할 수 있는 안정적인 바인딩 계층을 마련하였다. 이 중 \textbf{cuTile} DSL은 행렬 곱셈이나 어텐션 연산 등의 집약 연산을 타일(Tile) 단위로 기하학적으로 추상화하여, 레지스터 및 공유 메모리 캐싱 레이아웃을 Python 코딩 수준에서 제어할 수 있도록 돕는다. 이는 기존의 하드웨어 가상화 기반 실시간 GPU 커널 마이그레이션 및 상태 체크포인트-복구(C/R) 프레임워크인 CRAC 등과 달리, 소스 코드 수준에서 높은 이식성을 직접 확보하는 대안적 접근이다. 특히 CUDA 13.3부터는 Compute Capability 9.0 (Hopper) 아키텍처까지 균일하게 지원하므로, 로컬 데스크톱(RTX 40/50 시리즈)에서 작성된 어텐션 타일 커널이 클라우드 서버(H100/H200)로 추가적인 수정이나 재작성 공정 없이 투명하게 컴파일 및 포팅되는 강력한 하드웨어 호환성(Hardware Portability)을 확보할 수 있다.

\subsection{시스템 최적화 파트너로서의 코딩 에이전트}
기존의 LLM 코드 생성 및 자동 최적화 연구들은 정적으로 파라미터를 탐색하는 단순 그리드 서치(Grid Search)나 신택스 에러를 수정하는 자율 루프 시스템 모델에 치중해 왔다. 최근에는 하드웨어 프로파일링 피드백을 루프에 주입하거나 멀티 에이전트 기반으로 커널 최적화를 자동화하려는 Astra 및 KernelAgent 와 같은 에이전트 시스템들이 제안되고 있다. 그러나 하드웨어 수준의 커널 최적화는 컴파일러 메타데이터 해석, 레지스터 스필링 감지, 호스트-디바이스 메모리 할당 특성 등 다차원적 지식의 융합이 필요하다. 본 연구에서는 자율 구동되는 루프 시스템을 직접 개발하는 대신, 개발자와 긴밀히 상호작용하는 \textbf{코딩 에이전트(Coding Agent)}를 설계 및 미시적/거시적 최적화의 피드백 파트너로 활용함으로써 복잡한 GPU 시스템 튜닝을 인간의 가이드 하에 가속하는 '컴파운드 시스템 엔지니어링(Compound System Engineering)' 방법론을 구현하였다.


\section{인지 컴파일러 및 과학실험 기반 3계층 공동 설계 방법론}
단일 GPU의 극단적인 하드웨어 자원 제약 하에서 LLM 추론 엔진 커널을 효과적으로 마이그레이션 및 최적화하기 위해서는, 고수준 아키텍처 모델과 저수준 물리 연산 간의 긴밀한 상호작용이 필수적이다. 본 연구는 시스템 개발 과정의 과학적 엄밀성과 탐색 효율성을 보장하기 위해, 커널 엔지니어링을 \textbf{통제된 과학실험(Controlled Scientific Experiment)}으로 모델링하고, 인간 개발자와 AI 코딩 에이전트가 협력하여 시스템 모델을 스스로 진화시키는 \textbf{인지 컴파일러(Cognitive Compiler)} 방법론을 제안한다.

\subsection{프레임워크 개요 및 과학적 방법론 도입}
전통적인 컴파일러가 정적이고 고정된 규칙에 기초해 소스 코드를 기계어로 매핑하는 것과 달리, 본 연구에서 제안하는 \textbf{인지 컴파일러}는 실험 결과를 지식으로 동적 흡수하여 컴파일 제약 규칙을 점진적으로 갱신한다. 이는 GPU 커널 튜닝의 복잡한 설계 공간 속에서 시행착오(Trial-and-Error)를 소거하고 시스템 지식을 체계적으로 누적하기 위한 과학적 방법론 사이클에 기초한다.

\begin{table}[h]
\centering
\caption{커널 엔지니어링의 과학적 방법론 사이클 매핑}
\resizebox{\columnwidth}{!}{%
\begin{tabular}{|l|l|}
\hline
\textbf{과학적 방법론 단계} & \textbf{프레임워크 대응 요소 및 실행 메커니즘} \\ \hline
가설 수립 (Hypothesis) & Architect가 시스템 모델과 전역 지식 베이스를 바탕으로 최적화 방향 및 탐색 범위 정의 \\ \hline
실험 설계 (Experiment Design) & TechLead가 구체적인 코드 변형 패치 제안 및 컴파일 옵션 명세화 \\ \hline
실험 수행 (Execution) & Executor가 타겟 GPU 머신 상에서 배포된 패치를 빌드/실행하고 텔레메트리 데이터 계측 \\ \hline
데이터 분석 (Analysis) & TechLead가 수집된 실행 로그 및 Nsight 물리 메트릭을 분석하여 실험 보고서 작성 \\ \hline
성찰 및 지식 갱신 (Reflection) & Architect가 인과관계를 종합적으로 평가하고 획득한 한계 지식을 지식 베이스에 영속화 \\ \hline
\end{tabular}%
}
\end{table}

\subsection{3계층 위계 구조 및 통제 경계}
본 프레임워크는 정보의 추상화 수준, 실행 권한 및 하드웨어 접근 수준에 따라 \textbf{Architect(상위)}, \textbf{TechLead(중위)}, \textbf{Executor(하위)}의 3계층 위계로 분할되어 작용한다. 이는 과학실험실 팀의 역할 분담을 모방하여 각 에이전트와 인간의 행동 통제 영역을 엄격하게 규정한다.

\begin{itemize}
    \item \textbf{상위 계층: Architect (인간 개발자 + 클라우드 LLM)}: 시스템 궁극 의도(Intent) 정의, 전역 제약 수립, 지식 최종 승인을 담당한다. 타겟 머신의 로컬 코드 실행에 직접 개입하지 않으며, 거시적 구조 설계와 추상화된 물리 데이터 해석에만 집중한다.
    \item \textbf{중위 계층: TechLead-Agent (Antigravity CLI 코딩 에이전트)}: 지식베이스 규칙과 가설을 소스코드 패치 및 튜닝 탐색 범위로 번역(Translation)한다. 아키텍처 불변성과 컴파일러 제약 조건의 일관성 수호에 전념하며, 타겟 머신에서 코드를 직접 컴파일/실행하지 않는다.
    \item \textbf{하위 계층: Executor (인간 개발자 + 로컬 실행 환경 + Nsight)}: 패치를 물리적으로 로드하여 GPU 상에서 JIT 컴파일 및 론칭, 실행 에러에 대한 로컬 자율 디버깅 수행, 프로파일 텔레메트리 계측을 수행한다. 전역적 아키텍처 설계에는 관여하지 않으며, 오직 로컬 오류 극복과 정량 데이터 수집에 한정된다.
\end{itemize}


\section{micro-vllm 아키텍처 및 핵심 최적화}
본 장에서는 Windows OS와 단일 소비자용 GPU 환경에 맞추어 마이그레이션하고 최적화한 micro-vllm 서빙 엔진의 핵심 구조적 개선 사항을 기술한다.

\subsection{Gloo 기반 분산 통신 백엔드 이식}
vLLM과 같은 프로덕션 서빙 프레임워크는 대규모 모델의 분산 연산을 위해 Linux 계열에서만 작동하는 NCCL 백엔드를 기본적으로 사용한다. Windows 네이티브 환경에서 분산 텐서 연산의 컴파일 호환성을 보장하기 위해, 우리는 PyTorch Distributed의 \textbf{Gloo 백엔드}를 다중 디바이스 분산 스택에 이식하였다. 이를 통해 분산 통신 계층의 대대적인 코드 수정이나 이진 컴파일(Binary compilation) 충돌 없이, 텐서 병렬 선형 레이어(Linear Layer)의 논리 구조를 Windows 네이티브 실행 환경으로 완전히 투명하게 마이그레이션할 수 있었다.

\subsection{벡터화된 논리-물리 블록 캐시 재구성}
실제 서비스 루프 내에서 prefix-cached prefill 연산(기존에 보존된 역대 프롬프트 이력 캐시와 신규 사용자 입력을 융합하는 단계)을 수행할 때, 새롭게 인가된 Key, Value 텐서 모양(Shape)과 메모리 풀 내부의 논리적 시퀀스 크기 간의 불일치로 인한 런타임 셰이프 오류(Runtime shape mismatch)가 빈번히 유발된다.
우리는 이 셰이프 매칭을 CPU 연산 루프를 사용하지 않고 GPU 상에서 병렬적으로 처리하기 위해, \textbf{벡터화된 논리-물리 캐시 재구성 레이어}를 제안하고 적용하였다. context 길이 seqlen\_k에 해당하는 메모리 오프셋 주소를 텐서 인덱스 연산으로 변환하여 물리 주소로 고속 매핑한다. 이 벡터화된 배치 재배치 로직은 GPU 오버헤드를 sub-millisecond 단위로 격하시켜 호스트 병목을 차단하고, dynamic continuous batching 상황에서의 형상 무결성을 완벽히 보장한다.

\subsection{단일 커널 최적화 (Tiled GEMM \& FMHA 64x64)}
FMHA(Fused Multi-Head Attention) 커널 설계 시, 메모리 계층 간의 IO 비용을 축소하기 위해 중간 연산 데이터인 attention score를 공유 메모리 및 레지스터 영역에 상주시켰다. online softmax 알고리즘을 타일 레벨에 주입하여 수치 안정성을 취하고, 디코딩과 prefill 단계의 효율성을 도모하는 causal masking을 cuTile 스레드-블록 연산에 통합하였다.


\section{호스트-디바이스 인터페이스 최적화 및 자원 격리}
\subsection{JIT 컴파일 스톰과 메모리 할당자 스래싱 딜레마}
JIT 컴파일 기반 Python GPU DSL인 cuTile은 입력으로 전달되는 텐서 형상의 기하학적 변화에 반응하여 실시간으로 PTX 바이너리를 즉시 재빌드한다. continuous batching 서빙 환경 하에서는 동적으로 매 디코드 스텝마다 배치 크기와 시퀀스 길이가 변화하므로 JIT 컴파일이 무수히 유발되는 JIT Compilation Storm이 심각한 성능 지연을 만든다.
이러한 JIT 컴파일의 발생 빈도를 억제하기 위해, 우리는 텐서 형상의 차원을 256의 배수로 강제 정렬하는 \textbf{dynamic padding} 기법을 튜닝 옵션으로 적용해 보았다. 그러나 벤치마크 결과는 예상을 완전히 벗어나 성능이 급격히 저하되었다.
JIT 컴파일 엔진은 컴파일이 완료된 커널을 내부 해시 테이블(Hash Table)에 적재하여 캐싱하므로, 동일한 Shape이 유입될 때 컴파일 없이 수 $\mu$s 내에 끝나는 해시 룩업 지연만을 유발한다. 반면, dynamic padding을 우회 통과시키기 위해 매 디코드 스텝마다 가변 텐서를 256의 배수로 정렬하여 대형 패딩 텐서를 매번 새로 할당(New Allocation)할 경우, PyTorch의 캐싱 할당자는 내부 이중 연결 리스트 스캔을 반복 수행하여 호스트 CPU 연산 병목을 유발한다. 나아가, Eager-mode LLM Serving 하에서의 \textbf{암묵적 동기화(Implicit Synchronization)} 병목이 추가로 규명되어 실제 GPU 상의 연산 스케줄링을 연쇄적으로 방해하는 간섭 부작용이 실증되었다.

\subsection{CUDA Green Context 기반 SM 자원 분할 및 캐시 지역성 최적화}
우리는 CUDA Python 1.0의 최신 하드웨어 분할 API인 \textbf{Green Contexts}를 적용하여 RTX 5070 GPU의 전체 48개 SM 자원을 물리적으로 강제 파티셔닝하였다.
\begin{itemize}
    \item \textbf{Prefill Context (32 SMs)}: 막대한 부동소수점 연산 처리가 동반되는 프롬프트 분석용 SM 그룹.
    \item \textbf{Decode Context (16 SMs)}: 실시간 응답 지연을 방어하기 위한 1토큰씩의 생성용 전용 SM 그룹.
\end{itemize}
이러한 물리적 SM 구획화는 GPU 하드웨어 레벨에서 \textbf{L2 캐시 스레시(L2 Cache Thrashing) 및 데이터 방출(Eviction)}을 억제하는 미시적 메커니즘으로 작동하여, 디코드 연산 시 데이터 로드가 VRAM 버스로 우회하지 않고 L2 캐시 히트(L2 Cache Hit)로 다이렉트 처리되는 미시적 캐시 로컬리티 향상을 완성한다.


\section{실험 및 평가}
\subsection{실험 환경 및 시나리오}
Windows 11 OS 환경에서 NVIDIA GeForce RTX 5070 GPU(48 SM, 12GB VRAM, CUDA 13.3) 및 Qwen2.5-3B-Instruct 모델을 사용해 평가를 실시하였다. 동시 사용자 수는 256명으로 고정하고, 입력 및 출력의 크기를 100에서 1024 토큰 범위 내에서 동적으로 조율하는 continuous batching 상황을 재현하였다.

\subsection{패딩 여부에 따른 할당 오버헤드 실증 (Eager vs Padded)}
실험 측정 결과, dynamic padding 설정은 Eager 설정에 비해 전체 시스템 Throughput이 67.0\% 하락하는 최악의 병목을 유발하였다.

\begin{table}[h]
\centering
\caption{셰이프 패딩 적용 유무에 따른 E2E 서빙 효율성 비교}
\resizebox{\columnwidth}{!}{%
\begin{tabular}{|l|c|c|c|c|}
\hline
\textbf{시스템 아키텍처 설정} & \textbf{처리 총 토큰 수 (개)} & \textbf{총 소요 시간 (초)} & \textbf{생성 처리량 (tok/s)} & \textbf{성능 대조 비율 (\%)} \\ \hline
\textbf{Eager cuTile (본 연구)} & 133,966 & 284.54 & \textbf{470.82} & 기준치 (100.0\%) \\ \hline
\textbf{Padded cuTile} & 133,966 & 863.39 & 155.16 & -67.0\% (성능 하락) \\ \hline
\end{tabular}%
}
\end{table}

\subsection{Green Contexts 적용에 따른 성능 비교}
Green Contexts를 활성화하여 SM 자원을 Prefill 32개, Decode 16개로 분할 구획한 타겟 설정은 베이스라인 대비 \textbf{중간값 Decode 지연시간(P50 ITL)을 14.0\% 단축}시키는 성과를 보였다. 더불어 전체 토큰 처리량(Throughput) 역시 \textbf{10.6\% 증가}하는 고무적인 성능 우위를 보였다. 

\begin{table}[h]
\centering
\caption{Qwen2.5-3B-Instruct 기반 Baseline vs Green Contexts 성능 비교}
\resizebox{\columnwidth}{!}{%
\begin{tabular}{|l|c|c|c|}
\hline
\textbf{평가 성능 지표} & \textbf{Baseline (Green OFF)} & \textbf{Target (Green ON)} & \textbf{개선 비율 (\%)} \\ \hline
\textbf{TTFT (Prefill Latency)} & 244.60 ms & 243.42 ms & -0.5\% (유지) \\ \hline
\textbf{Decode P50 ITL} & 48.63 ms & \textbf{41.83 ms} & -14.0\% (단축) \\ \hline
\textbf{Decode P99 ITL} & 85.31 ms & 84.34 ms & -1.1\% (안정화) \\ \hline
\textbf{Total Throughput} & 415.50 tok/s & \textbf{459.73 tok/s} & +10.6\% (향상) \\ \hline
\end{tabular}%
}
\end{table}


\section{결론 및 향후 연구}
본 논문은 Windows OS 환경에서 cuTile DSL과 CUDA Python 1.0을 결합한 독창적인 고성능 LLM 추론 엔진인 micro-vllm을 구현하고 E2E 수준의 병목 요인을 실증 규명하였다. 특히 코딩 에이전트와의 밀접한 협업 설계 및 미시 프로파일링 피드백을 통해 64x64 타일링과 벡터화된 논리-물리 주소 변환 레이어를 구축하였으며, eager 모드 서빙 시 JIT 방지용 패딩 기법이 가혹한 메모리 할당자 스래싱을 유발하여 Throughput을 67.0\% 저하시키는 역효과가 있음을 정량적 벤치마크를 통해 실증적으로 규명하였다. 아울러 Green Contexts 기반 SM 파티셔닝 튜닝을 통해 Decode P50 ITL을 14.0\% 단축하는 시스템 개선을 실현하였다.
향후 연구로는 엔진 루프를 비동기 멀티스레드(Asynchronous Multi-threaded) 아키텍처로 전면 이원화하여, 서로 다른 GPU 스트림과 물리 SM 파티션 상에서 Prefill과 Decode가 오버랩(Overlap)되어 동작하는 완전한 E2E 자원 격리를 구현해 나갈 계획이다.

\section*{감사의 글}
본 연구의 코드 최적화 및 실험 설계 과정에서 AI 코딩 에이전트(Google Antigravity CLI)를 보조 도구로 활용하였으며, 모든 최종 결정과 검증은 인간 저자가 수행하였다.

% 부록 및 참고문헌 추가 시 활용 가능
%\bibliographystyle{ieeetr}
%\bibliography{references}

\newpage
\appendix
\section{부록 A: cuTile 프로그래밍 모델 핵심 개념}
NVIDIA가 도입한 cuTile DSL은 전통적인 CUDA C++의 극단적인 저레벨 복잡성과 PyTorch의 고수준 추상화 사이의 간극을 메우는 타일 기반 프로그래밍 패러다임이다.
\begin{itemize}
    \item \textbf{타일 추상화 (Tile Representation)}: 모든 행렬 및 벡터 데이터는 고정된 크기의 기하학적 블록인 타일(Tile)로 정의되며, 컴파일러가 스레드 레이아웃 배치를 전 자동화한다.
    \item \textbf{MMA 하드웨어 제어}: Tensor Core 하드웨어가 직접 처리하는 고속 행렬 연산 지시어인 MMA를 \texttt{ct.mma}로 제공하여 연산 능력을 극대화한다.
    \item \textbf{협력적 메모리 로드 및 스토어}: VRAM과 레지스터 간 복사는 \texttt{ct.load} 및 \texttt{ct.store} 명령어를 통해 블록 내 스레드들이 공유 메모리를 매개로 공동 집행한다.
    \item \textbf{Hopper 아키텍처 이식성}: 단일 소비자용 RTX GPU에서 컴파일되는 동일한 Python 코드가 엔터프라이즈 전용인 H100/H200 머신으로 진입 시, 100\% 투명하게 컴파일을 승계할 수 있다.
\end{itemize}

\section{부록 B: LLM 추론 동작 원리 핵심 개념}
\begin{itemize}
    \item \textbf{Prefill 단계}: 사용자가 주입한 입력 텍스트 전체를 동시에 파싱하며, 연산 성능 제한적(Compute-bound) 성격을 가진다.
    \item \textbf{Decode 단계}: 오직 1개의 다음 토큰을 순차 예측하며, 메모리 대역폭 제한적(Memory-bandwidth bound) 특성이 병목을 지배한다.
    \item \textbf{Paged KV Cache}: 가상 메모리 기법을 빌려 VRAM의 캐시 영역을 일정한 블록(Block) 단위로 물리 페이지로 분할하고 단편화를 차단하는 핵심 기술이다.
    \item \textbf{Continuous Batching (연속 배치 처리)}: 각 시퀀스가 개별적으로 종료되면 즉시 새로운 요청을 동적으로 교체 투입하여 유연성을 제공하는 기법이다.
\end{itemize}

\end{document}
